% ===== 原文 pp.81--84 =====
% [原文 p.81]
\begin{exercise}{7.1}
回忆带四个穿刺点的球面的 $\fsu(2)$ class S 理论的 SW 曲线 (4.13)：$x^2=u_2(z)$，其中
\begin{equation}
    u_2(z)=\frac{\frac{q}{z}m_1^2+\frac{q(q-1)}{z-q}m_2^2+\frac{z-q}{z-1}m_3^2+zm_4^2-u}{z(z-q)(z-1)}.
    \tag{7.20}
\end{equation}
该理论有一个描述，即规范耦合为 $\tau=(\log q)/(2\pi\ii)$、带 $N_f=4$ 个味（质量分别为 $m_1\pm m_2$ 与 $m_3\pm m_4$）的 $SU(2)$ SQCD。

1. 退耦一个超多重态：取 $m_1+m_2\to\infty$，保持 $m_1-m_2$、$m_3\pm m_4$ 以及 $\Lambda=q(m_1+m_2)$ 固定。你应当得到 $u_2=P(z)/(z^4(z-1)^2)$，其中 $P$ 是某个四次多项式。

2. 用两种方式退耦第二个超多重态。第一，取 $m_1-m_2\to\infty$，保持 $\Lambda'^2=\Lambda(m_1-m_2)$ 及其他质量固定。第二，改为取 $m_3+m_4\to\infty$，保持 $\tilde z=z(m_3+m_4)$、$\tilde x=x/(m_3+m_4)$ 以及 $\Lambda'^2=\Lambda(m_3+m_4)$ 和其他质量固定。把一条 SW 曲线映射到另一条，并检验两者的 SW 微分 $\lambda$ 之差是不重要的（留数是质量，没有 Coulomb 分支依赖）。

3. 退耦第三个与第四个超多重态，并重新标度 $z\to z\Lambda^2$，得到众所周知的纯 $SU(2)$ SYM 曲线：$z^2x^2=u+\Lambda^2(z+1/z)$，其中 $\lambda=x\dd z$。
\end{exercise}
% [Tutor Note: p.81 习题 7.1 第 1 步可自己验算：令 $M=m_1+m_2\to\infty$、$q=\Lambda/M$，则 $\frac{q}{z}m_1^2\simeq\frac{\Lambda M}{4z}$、$\frac{q(q-1)}{z-q}m_2^2\simeq-\frac{\Lambda M}{4z}$，领头 $M^2$ 项相消，只剩由 $\Lambda$ 控制的有限项，于是 $u_2=P(z)/(z^4(z-1)^2)$：$z=0$ 处的 4 阶极点是两个驯服穿刺点（$m_1$ 在 $z=0$、$m_2$ 在 $z=q$）相撞后的野生穿刺点，阶数 $4=2\times2$ 与两个驯服点相撞一致；$z=1$ 处 $m_3$ 的 2 阶极点与 $z=\infty$ 处（$zm_4^2$ 项）的重数都保持不变。]

\para{CFT 一侧。} 在 2d CFT 一侧，产生野生穿刺点的那些极限对应于让主顶点算子 $\hat V_\alpha$ 相撞。$n\ge2$ 个顶点算子相撞产生记为 $I_{n-1}$ 的非正规算子（irregular operators），它们依赖于原来的 $n$ 个动量（经过适当标度）。它们与能动张量的 Ward 恒等式所涉及的极点，其阶数与规范理论一侧在同一次碰撞中出现的 $\phi_2$ 极点阶数相同，都是 $2n$。这一吻合与“$\phi_2$ 可理解为 $T$ 的半经典极限”这一事实一致：
\begin{equation}
    T(z)I_{n-1}(0)=\sum_{k=n-1}^{2n-2}\frac{\Lambda_k}{z^{k+2}}I_{n-1}(0)+\cdots
    \;\stackrel{r\to\infty}{\simeq}\;
    r^2\phi_2(z)I_{n-1}(0)+\ldots.
    \tag{7.21}
\end{equation}
由态-算子对应（state-operator correspondence），这些算子给出 Virasoro 代数的相干态（coherent states）[94]（又称 Whittaker 向量或 Gaiotto 态），以及它们更一般的推广------称为非正规态（irregular states，有时也叫 Bonelli--Maruyoshi--Tanzini（BMT）态）[366,367]。

\para{野生穿刺点与 AD 理论。} 从 $SU(2)$ $N_f=4$ SQCD 或 $\cN=2^*$ SYM 出发退耦超多重态、以及它在 AGT 对应中的效应，在 [94,368--375] 中研究过。Virasoro 代数更高阶的非正规态（BMT 态）在 [366,367,376--380] 中被研究，推广到 W 代数则在 [99,378,381--383]。在 CFT 一侧，主算子的碰撞与非正规态的直接定义由 Rim 及其合作者 [377,379,380,384--391] 以及其他人 [68,392--401] 给出。
% [Tutor Note: p.81 原文此处把 correspondence 误拼为 "corespondence"（"its effect in the AGT corespondence"），译文按正确词义翻译；不影响内容，仅记录以便校勘。]

% [原文 p.82]
\para{进一步的方向。} 一些 AD 理论是在 Lagrangian 规范理论的 Coulomb 分支的特定点上、以及作为 S 对偶描述中出现的理论而被发现的 [2,169,363--365,402]。各种 AD 理论及相关课题的 class S 构造见 [6,8,403--410]，以及 Xie 与其合作者的工作 [411--423]。人们并不指望 class S 理论穷尽所有可能的 4d $\cN=2$ 理论（例如见 [424]）。除了 Lagrangian 场论的分类 [425,426]、以及按 Coulomb 分支几何对理论进行分类的纲领 [427--435] 之外，人们还探索了 4d $\cN=2$ 理论的其他构造 [406,414,416,417,419,436--447]。

一个有趣的工具可以在 class S 理论没有 Lagrangian 描述时仍然检验 AGT 对应，那就是计算中心荷与反常。Toda CFT（及其推广）的中心荷，在 [19,22,448,449] 中与 6d (2,0) 理论 $\cX(\fg)$ 的反常 8-形式（anomaly 8-form）向 2d 的约化相匹配。若改为在 $C$ 上约化，就给出 4d class S 理论的 $a$ 与 $c$ 共形反常 [100,117,120,450--453]。

\section{各种维数的算子}

Wilson [454] 与 't Hooft [455] 圈算子，以及把二者结合起来的双子（dyonic）圈 [456]，在研究 4d 规范理论的相（phases）时起重要作用。面算子（surface operators）研究得较少，却非常丰富；例如某些面算子的模空间就是 class S 理论的 UV 曲线。最后，畴壁（domain walls）描述两个 4d 理论之间的界面，或者边界现象。在 AGT 对应中，这类算子的一个巨大来源是 6d (2,0) 理论 $\cX(\fg)$ 的半 BPS 余维 2 与余维 4 缺陷。另一个来源是把 6d 设置作轨形化（orbifold），而轨形化群必须尊重定向，因为 $\cX(\fg)$ 是手征理论。

这些缺陷可以按相对于乘积时空 $M_4\times C$ 的各种定向插入到 AGT 对应中。我们已经详细讨论过余维 2 缺陷插入在 $C$ 中一点、并包裹整个 4d 时空的情形：这些实际上就是贯穿本综述的穿刺点与扭算子，尤其是在第 7 节。

\begin{keybox}{规范理论与 CFT 算子的主要例子}
\begin{itemize}[leftmargin=1.4em,itemsep=1pt,topsep=1pt]
  \item 余维 2 情形：与一个 tinkertoy 耦合对应于插入一个顶点算子；对称性破缺的壁对应于 Verlinde 圈；Gukov--Witten 面算子会改变 Toda CFT。
  \item 余维 4 情形：涡旋弦面算子对应于退化顶点算子；Wilson--'t Hooft 圈对应于退化的 Verlinde 圈；$S^4_b$ 极点处的轨形奇点会改变 Toda CFT。
\end{itemize}
\end{keybox}

关于已被研究过的全部可能性及参考文献，我们请读者参见引言中的 Table 1。这里我们按 4d 一侧维数递增的顺序组织讨论：先讨论 8.1 小节中的点状算子，然后是 8.2 小节中的线与圈算子（见综述 [457]）、8.3 小节中的 2d 算子（见综述 [458]），以及 8.4 小节中的 3d 壁与界面。

% [原文 p.83]
\subsection{4d 中的局域算子}

$\cX(\fg)$ 的余维 4 算子由 $\fg$ 的表示标记，而把这样的算子包裹整个 $C$ 的效应尚未被完全理解。事实上，对我们下面描述的大多数几何，主要想法在某些理论中（例如 $\fg=\fsu(2)$）已经理解，但并非在完全一般的情形。

\para{Coulomb 分支算子。} 定义 SW 曲线的 $k$ 阶全纯微分 $\phi_k(z)=u_k(z)\dd z^k$ 可以由 Toda CFT 关联函数的经典极限算出。在这个极限中（$S^4_b$ 的半径很大，或 $\epsilon_1,\epsilon_2\to0$，等价地 2d CFT 的共形维数很大），二次微分 $u_2$ 由能动张量的插入给出，更一般地 $u_k$ 由自旋 $k$ 的流 $W_k$ 给出：\footnote{数值系数依赖于约定。}
\begin{equation}
    u_k(z)\propto\frac{\langle W_k(z)\hat V_{\mu_1}\ldots\hat V_{\mu_n}\rangle}{\langle\hat V_{\mu_1}\ldots\hat V_{\mu_n}\rangle}
    \qquad\text{as }\epsilon_1,\epsilon_2\to0.
    \tag{8.1}
\end{equation}
对 $\fsu(2)$ 见原始 AGT 论文 [5]，或更显式的 [459] 等。

现在考虑 $C$ 的一个裤子分解，以及把 $\mathsf{T}(\fg,C,D)$ 相应地描述为一组 tinkertoy 与规范化某些对称性的矢量多重态。设 $\varphi$ 是其中一个矢量多重态（对应一根管子）中的标量，并考虑 A 型情形下诸如 $\Tr\varphi^l$ 的规范不变式，更一般地 $\fg$ 的所有 Casimir。经典地，它们作为微分 $\phi_k$ 的系数出现，因而可以作为 $\phi_k$ 的某些带权积分被取回。回到一般的 $\epsilon_1,\epsilon_2$，4d 一侧的算子 $\Tr\varphi^l$ 对应于流 $\tilde W_l$ 的一个适当的带权积分 [19]。\footnote{自旋 $l$ 的流 $\tilde W_l$ 是 $W_k$ 的多项式，以计入 Casimir $\Tr\varphi^l$ 与特征多项式 $\det(x-\varphi)$ 的系数之间的差别。} 例如，插入 $\Tr\varphi^2$ 相当于对 $Z_{S^4_b}$ 关于规范耦合求导数 [460--462]，而规范耦合对应于 $C$ 的形状，这确实翻译成全纯能动张量 $T=\tilde W_2$ 的积分插入。
% [Tutor Note: p.83 脚注 62 的内容可以具体验证：对 $\fsu(N)$，$\det(x-\varphi)$ 中 $x^{N-l}$ 的系数是初等对称多项式 $e_l(\varphi)$（即 $\wedge^l\varphi$ 的迹），而 $\Tr\varphi^l$ 是幂和 $p_l$。由牛顿恒等式 $p_l-e_1p_{l-1}+e_2p_{l-2}-\cdots+(-1)^l l\,e_l=0$ 可知 $l\ge3$ 时 $e_l$ 含有 $p_l$ 之外的多项式项（例如 $e_2=\frac12(p_1^2-p_2)$），故 $W_k$ 的多项式组合才能给出 $\Tr\varphi^l$；$l=1,2$ 时二者只差常数倍。这也是 $\tilde W_l$ 用 $\tilde{}$ 区分于 Toda 的 $W_k$ 的原因。]

带（乘积的）$\Tr\varphi^j$ 插在一个极点、$\Tr\bar\varphi^k$ 插在另一个极点的 $S^4_b$ 上的关联函数，可以用超对称局域化计算，虽然这些算子会使瞬子计数复杂化。通过共形变换，圆的情形 $b=1$ 给出平坦空间上恰含一个反手征算子的关联函数的结果 [192,463--467]，这些结果为各种场论想法提供了细致的检验，例如复活（resurgence）[192,468]、大电荷展开（large-charge expansions）[469--476] 以及更多 [477]。这些特定的关联函数在对应的 2d CFT 一侧尚未被研究过。

\para{轨形 $\C^2/\Z_M$。} 接下来我们考虑另一个操作，其效应是让 $\R^4$ 内的一点（或 $S^4_b$ 的两个极点）变奇异：用群 $\Z_M$ 作轨形化，它在 $\C^2$ 上作用为 $(\exp(2\pi\ii/M),\exp(-2\pi\ii/M))$。超对称局域化仍然可行：只需把所有模限制为 $\Z_M$ 不变的模，并把瞬子计数限制为 $\Z_M$ 不变的

% [原文 p.84]
瞬子计数。这似乎对应于陪集 CFT
\begin{equation}
    \frac{\widehat{\fsu}(N)_k\times\widehat{\fsu}(N)_M}{\widehat{\fsu}(N)_{k+M}}\times\frac{\widehat{\fsu}(M)_N}{\widehat{\fu}(1)_{M-1}},
    \qquad k=-N-\frac{Mb^2}{1+b^2},
    \tag{8.2}
\end{equation}
其能级 $k$ 是分数，而 $M=1$ 时它会非平凡地约化为通常的 Toda CFT。$N=M=2$ 的情形本质上就是 super-Liouville CFT，在 [20,23--25,27,29,30,32,478,479] 中研究过，另见 [63] 中带一个面算子的情形。在 $\C^2/\Z_2$ 及其爆破（此时在 $U(1)$ 等距的两个不动点处各有一份瞬子）上的瞬子计数是相关联的 [230,233,480,481]。这导致把 super-Liouville CFT 分解为 Liouville 与类时 Liouville（time-like Liouville）两部分 [317,482--484]。一般的 $N,M$ 推广在 [22,26,28,31,34] 中被部分完成，而陪集 (8.2) 在 [485--487] 中得到进一步研究。另一个视角是把 $\R^4/\Z_M$ 实现为 $\R^4\times S^1$ 的维数约化，这对应于在 9.1 小节的 q-形变 AGT 对应中把 $q$ 取为单位根 [21,33,35--37]。
% [Tutor Note: p.84 (8.2) 的两点核对：(i) 能级 $k=-N-\frac{Mb^2}{1+b^2}$ 是分数，这正是 $M\ge2$ 时该陪集称为分数级（fractional level）WZW 模型的原因；$M=1$ 时分子分母的第一因子相消，只剩 $\widehat{\fsu}(N)_k\times\widehat{\fsu}(1)_N/\widehat{\fu}(1)_0$ 型结构，回到 Toda。(ii) $N=M=2$ 时分子为 $\widehat{\fsu}(2)_k\times\widehat{\fsu}(2)_2$、分母为 $\widehat{\fsu}(2)_{k+2}$，即 super-Liouville 的 $\widehat{\fsu}(2)$ 结构；这与把 $\C^2/\Z_2$ 上的瞬子按两个不动点分解、从而把 super-Liouville 拆成 Liouville 与类时 Liouville 的做法相对应。]

\subsection{线算子}

我们现在把 6d 理论的余维 4 算子放在 $L\times\gamma$ 上，其中 $L$ 是 $S^4_b$ 中由超对称性允许的圆周之一：$\{y_3=y_4=0,\;y_5=\mathrm{const}\}$ 或 $\{y_1=y_2=0,\;y_5=\mathrm{const}\}$，而 $\gamma$ 是 $C$ 中一条无自交的闭圈。经维数约化，该算子在 AGT 对应中插入圈算子：一个由 $\gamma$ 标记、置于 $L\subset S^4_b$ 上的圈算子（在配分函数中），以及一个 $\gamma$ 上的圈算子（在 Toda CFT 关联函数中）。这套设置在 $\fsu(2)$ 情形于 [38--43] 中研究过，对更一般的 $\fg$ 在 [45--52]，对这类算子的网络在 [44,51--54]，对线算子代数在 [343,488,489]，另见 [490,491]。线算子与圈算子在刻画 4d 规范理论的相 [454--456]、以及精细化规范群的全局结构 [492] 时起重要作用；对 class S 见 [38,50--52,493--500]。4d $\cN=2$ 规范理论中圈算子的精确期望值用超对称局域化计算，见 [10,501,502]，综述见 [457]，其中重要的微妙之处后来在 [503--508] 中得到澄清。关于 4d 理论中线缺陷可观测量的其他考虑包括 [509--515]。

\para{Wilson 圈算子。} 由于 $\gamma\subset C$ 无自交，我们可以沿它切开 $C$，得到带两个额外穿刺点（带有某些数据，比如 $D_1,D_2$）的（可能不连通的）曲面 $C'$。如 (1.2) 附近所讨论，相应的 class S 理论 $\mathsf{T}(\fg,C,D)$ 由对应于 $C'$ 的理论 $\mathsf{T}(\fg,C',\{D,D_1,D_2\})$ 通过规范化与 $D_1,D_2$ 相关联的味对称性 $F_1,F_2$ 的对角子群而得到，如 (1.2) 那样。这样，每条无自交的圈 $\gamma$ 在 $\mathsf{T}(\fg,C,D)$ 的某个描述中对应于一个规范群 $G_\gamma=(F_1\times F_2)_{\mathrm{diag}}$。

在 $G_\gamma$ 弱耦合的极限下，圈 $\gamma$ 包裹一根细管子。把 $\cX(\fg)$ 在这根管子上约化给出区间上的 5d $\cN=2$ SYM，而余维 4 缺陷给出一个线算子，具体地是一个 Wilson 圈，它额外地依赖于 $\fg$ 的表示 $R$ 的选择。事实上，这是看清余维 4 算子
